\section{AION Flow Intelligence Stack Catalogue and Canvas}
\label{sec:aion-flow-intelligence-stack-catalogue}

\subsection{Purpose and Product Boundary}

AION Flow extends the existing Glyph workflow canvas into a visual, provider-neutral
intelligence orchestration environment.  It does not create a second workflow product.
The existing \emph{Add a step} catalogue, canvas, node graph, autosave mechanism and full
Node Editor remain the canonical authoring surfaces.  A new \emph{Intelligence Stack}
category exposes replaceable models, harnesses, compute and evidence while the customer's
AION brain retains identity, memory, policy, authority, verified outcomes and learning.

The governing execution shape is:

\begin{quote}
AION admission $\rightarrow$ permitted context and evidence $\rightarrow$ specialist
intelligence $\rightarrow$ bounded deliberation $\rightarrow$ AION policy and approval
$\rightarrow$ capability execution $\rightarrow$ verification $\rightarrow$ receipt and
outcome learning.
\end{quote}

The AION boundary therefore surrounds the intelligence workflow.  A model is neither the
owner of the workflow nor an execution authority.  Its output is a proposal that must pass
evidence, policy, permission and verification controls.

\subsection{Catalogue Families}

The new catalogue presents eight searchable families:

\begin{enumerate}
  \item \textbf{AION Brain}: admission, active identity and space, scoped memory, Business
  Map context, purpose and receipt boundaries;
  \item \textbf{Models}: local models and customer-controlled private endpoints;
  \item \textbf{Harnesses}: signed instructions, schemas, procedures, examples and declared
  tools that can be versioned independently from a model;
  \item \textbf{Compute}: the mother brain, another AION node, a private server or a
  customer-owned cloud target;
  \item \textbf{Evidence}: permissioned retrieval and deterministic schema, calculation,
  policy and business-rule validation;
  \item \textbf{Deliberation}: independent criticism and evidence-based reconciliation with
  explicit iteration, time and cost ceilings;
  \item \textbf{Governance}: disclosure, redaction, residency, encryption and private human
  approval; and
  \item \textbf{Action and Verification}: idempotent governed capabilities, observation,
  verification and normalized route receipts.
\end{enumerate}

The initial catalogue deliberately supplies structural nodes rather than embedding a single
provider's product vocabulary.  For example, a \emph{Private model endpoint} uses an opaque
endpoint binding and immutable model manifest; provider-specific implementations can later
bind to that contract without changing the graph.

\subsection{Pre-insertion Disclosure}

Each Intelligence Stack card declares its execution location, credential requirement, data
classes, indicative cost, risk, maturity, modality, domain, licence, approval requirement and
whether it can cause an external effect.  Local, customer-cloud and external destinations are
visually differentiated.  A compact boundary legend explains that any route leaving the
customer boundary must declare its destination, data classification, encryption and cost before
execution.

The module search operates across names, descriptions, node families, locations, modalities,
domains, licences and maturity metadata.  The eight families can also be selected as compact
filters.  Search controls, family tabs and module rows expose accessible labels and keyboard
focus behaviour.

\subsection{Schema-driven Node Configuration}

Double-clicking an Intelligence Stack node continues to open the existing full Node Editor.
The editor retains its Input Schema, Input Table, Input JSON, Parameters, Settings, Output
Schema, Output Table and Output JSON views.  An additional schema-driven panel is rendered for
the selected intelligence module.  This panel provides:

\begin{itemize}
  \item typed string, number, integer, Boolean, enumeration and list controls;
  \item safe defaults, minimum and maximum values, required-field indicators and contextual help;
  \item a non-executing configuration test that reports missing required settings;
  \item a reset operation that restores the module's safe defaults; and
  \item a concise preflight summary of location, credentials, risk, cost, approval and possible
  external effect.
\end{itemize}

Configuration testing validates structure only.  It never invokes a model, connector or target
system.  Consequential nodes keep the existing execution control disabled until the compiled
flow proves policy, authority and an exact, unexpired approval.

\subsection{Credential and Browser-state Safety}

Raw API keys, passwords, provider tokens, private keys and secret-shaped values are rejected by
the node editor.  A node stores only an opaque reference of the form
\texttt{vault://binding-name}.  The referenced secret remains in the encrypted mother-brain
vault and is resolved at execution time under the active identity, purpose and policy.  The raw
secret is excluded from the canvas graph, browser state, autosave, logs, previews, receipts and
workflow exports.

This indirection permits credential rotation, revocation and provider replacement without
editing every workflow.  Exported graphs remain inspectable and portable but cannot carry a
customer's usable credentials.

\subsection{Compatibility with Existing Glyph Workflows}

Intelligence nodes are inserted through the same graph mutation path as existing Flow Control,
Tools, Text Parser, Apps, Approval and Custom Logic nodes.  Their module identity, safe
configuration schema and disclosure metadata are stored inside the existing node configuration
object.  Existing autosave, undo and redo, zoom, copy and paste, edge routing and node-editor
behaviour remain in place.  No existing workflow is migrated or executed merely because the new
catalogue is available.

\subsection{Implemented Initial Nodes}

The first implementation includes AION admission, scoped Business Map context, local model,
private model endpoint, versioned harness, configurable compute target, evidence retrieval,
deterministic validation, independent critic, evidence-based reconciliation, disclosure and
redaction gate, private human approval, governed capability, and verify-and-receipt nodes.

The visible product now supports safe graph composition and configuration.  Live provider
execution remains intentionally outside this stage.  Subsequent stages add signed model and
harness registries, qualified compute adapters, richer evidence acquisition, deliberation,
governance visualization, dry-run compilation and finally permissioned execution.

\subsection{Verification Status}

The JavaScript bundle passes syntax validation.  Automated interface contracts verify the new
catalogue, all eight families, module insertion metadata, search and filters, disclosure badges,
schema-generated configuration, raw-secret rejection, opaque vault references and the
consequential-action execution lock.  Existing workflow picker and review-panel contract tests
also pass.  The live Pilot television and AION learning processes do not require a restart for
this documentation stage and were not interrupted.
